% Generated semantic source for AI Almanach v3.
% Edit v3 directly after initial generation; do not overwrite
% editorial changes by rerunning the converter casually.

% ==== International IT Law =============================================================================
\chapter{International IT Law}

\chapterlead{Which duties attach to an AI system, its data, and its operators?}{Classify}{Identify}{Document}{Reassess}{it-law}

\begin{v3aside}[Chapter Scope]
    This is a \textbf{signpost chapter} for AI practitioners — not a legal textbook.
    It maps the regulatory landscape you will encounter when building, training, or deploying AI systems,
    and points you to the right regulations and references.
    When legal questions become concrete, consult qualified legal counsel.
    Regulatory status in this chapter was checked on 26 August 2026; always
    recheck the official text, application date, jurisdiction, and guidance.
\end{v3aside}

% ---- 17.1 Legal Foundations for AI Practitioners ------------------------------------------------
\section{Legal Foundations for AI Practitioners}

\subsection{Key Areas of IT Law}

\begin{v3topic}{IT Law: Areas and AI Relevance}
    \begin{tblr}{|X[2]|X[3]|X[3]|}
        \hline
        \textbf{Area} & \textbf{Core Concern} & \textbf{Relevant for AI When\ldots} \\
        \hline
        Data Protection / Privacy & Lawful processing of personal data & Training on personal data; profiling; automated decisions \\
        \hline
        Intellectual Property & Copyright, patents, trade secrets & Training data sources; model outputs; software licensing \\
        \hline
        AI-Specific Regulation & Risk-based obligations for AI systems & Deploying AI in the EU; high-risk use cases (hiring, credit, medical) \\
        \hline
        Cybersecurity Law & Security obligations, incident reporting & Running AI infrastructure; critical-sector deployment \\
        \hline
        E-Commerce / Contracts & Platform liability, online contracts & Selling AI services; deploying via APIs or SaaS \\
        \hline
        Product Liability & Liability for defective products / AI & AI causing harm to individuals or property \\
        \hline
    \end{tblr}
\captionof{table}{Key Areas of IT Law: IT Law: Areas and AI Relevance.}
\label{tab:ch17-it-law-inline-01}
    \textsuperscript{\cite[][pp.~1--25]{lloydInformationTechnologyLaw2020}}

\end{v3topic}
\begin{v3topic}{International vs.\ EU Law}
    Digital systems are inherently cross-border, yet law is territorial.
    Key challenges: \emph{which law applies?} (jurisdiction), \emph{which court decides?} (forum), \emph{will judgments be recognised?}
    The EU has emerged as the dominant regulatory force in digital markets — the \emph{Brussels Effect} means EU rules often become de-facto global standards because multinationals find a single compliance approach more practical\textsuperscript{\cite[][pp.~3--18]{savinEUInternetLaw2017}}.

\end{v3topic}
\begin{v3topic}{Regulatory Bodies}
    \begin{tblr}{|X[2]|X[4]|}
        \hline
        \textbf{Body} & \textbf{Role} \\
        \hline
        European Data Protection Board (EDPB) & Harmonises GDPR interpretation across EU \\
        \hline
        National DPAs & Enforce GDPR (e.g.\ BfDI Germany, ICO UK, CNIL France) \\
        \hline
        EU AI Office & Oversees EU AI Act compliance; established 2024 \\
        \hline
        FTC (US) & Enforces against unfair/deceptive AI practices in the US \\
        \hline
    \end{tblr}
\captionof{table}{Key Areas of IT Law: Regulatory Bodies.}
\label{tab:ch17-it-law-inline-02}
\end{v3topic}

% ---- 17.2 Data Protection and Privacy ----------------------------------------------------------
\section{Data Protection and Privacy}

\subsection{GDPR Essentials}

\begin{v3topic}{GDPR Core Principles (Art.\ 5)}
    The \textbf{General Data Protection Regulation} (EU 2016/679) governs any processing of personal data of EU residents, regardless of where the processor is located\textsuperscript{\cite[][pp.~15--30]{itgovernanceprivacyteamEUGDPR2020}}.
\visualseparator
    \begin{itemize}
        \item \textbf{Lawfulness, fairness, transparency} — processing must have a legal basis
        \item \textbf{Purpose limitation} — collected for specified, explicit purposes only
        \item \textbf{Data minimisation} — only what is necessary
        \item \textbf{Accuracy} — kept accurate and up to date
        \item \textbf{Storage limitation} — not kept longer than necessary
        \item \textbf{Integrity and confidentiality} — appropriate security
        \item \textbf{Accountability} — controller must demonstrate compliance
    \end{itemize}

\end{v3topic}
\begin{v3topic}{Legal Bases for Processing (Art.\ 6)}
    Every processing activity requires one of six legal bases\textsuperscript{\cite[][pp.~31--50]{itgovernanceprivacyteamEUGDPR2020}}:
    \begin{tblr}{|X[2]|X[4]|}
        \hline
        \textbf{Basis} & \textbf{When to use} \\
        \hline
        Consent & User freely, specifically, informed consent \\
        \hline
        Contract & Processing necessary to perform a contract \\
        \hline
        Legal obligation & Required by law \\
        \hline
        Vital interests & Life-or-death situations \\
        \hline
        Public task & Official authority or public interest \\
        \hline
        Legitimate interests & Balance test; cannot override data subject rights \\
        \hline
    \end{tblr}
\captionof{table}{GDPR Essentials: Legal Bases for Processing (Art.\ 6).}
\label{tab:ch17-it-law-inline-03}
    For ML: \emph{legitimate interests} or \emph{consent} are most common, but each has risks.

\end{v3topic}
\begin{v3topic}{Data Subject Rights}
    GDPR grants individuals enforceable rights that AI systems must
    support\textsuperscript{\cite[][pp.~80--120]{itgovernanceprivacyteamEUGDPR2020}}:
    \par\smallskip\noindent
    \begin{tblr}{width=\linewidth,colspec={X[2]X[4]X[4]},hlines,vlines}
        \hline
        \textbf{Right} & \textbf{Meaning} & \textbf{AI Implication} \\
        \hline
        Access (Art.\ 15) & See what data is held & Models trained on personal data must allow data inventory \\
        \hline
        Rectification (Art.\ 16) & Correct inaccurate data & Correction must propagate to training sets where feasible \\
        \hline
        Erasure (Art.\ 17) & ``Right to be forgotten'' & May require retraining or machine unlearning \\
        \hline
        Portability (Art.\ 20) & Receive data in machine-readable format & APIs and data export pipelines required \\
        \hline
        Object (Art.\ 21) & Opt out of processing & Must honour objections to profiling \\
        \hline
    \end{tblr}
\captionof{table}{GDPR Essentials: Data Subject Rights.}
\label{tab:ch17-it-law-inline-04}
\end{v3topic}

\subsection{GDPR for AI/ML}

\begin{v3topic}{Automated Decision-Making (Art.\ 22)}
    Individuals have the right \emph{not to be subject to a decision based solely on automated processing} that produces \textbf{legal or similarly significant effects} — e.g.\ automated loan refusal, CV screening, parole assessment\textsuperscript{\cite[][pp.~140--155]{itgovernanceprivacyteamEUGDPR2020}}.
\visualseparator
    \begin{itemize}
        \item Exceptions: explicit consent, contractual necessity, or authorised by law
        \item Even where allowed: must provide \emph{meaningful information about the logic involved} and the right to \emph{human review}
        \item Practical implication: black-box models in high-stakes decisions are legally risky
    \end{itemize}

\end{v3topic}
\begin{v3topic}{Data Protection Impact Assessment (DPIA)}
    A \textbf{DPIA} (Art.\ 35) is mandatory before processing that is ``likely to result in a high risk'' to individuals — this includes large-scale profiling, systematic monitoring, and use of sensitive data\textsuperscript{\cite[][pp.~160--175]{itgovernanceprivacyteamEUGDPR2020}}.
    For AI: required when training on sensitive data (health, biometrics), deploying facial recognition, or building behavioural profiling systems.
\visualseparator
    DPIA documents: purpose, necessity, proportionality, risks, and mitigations.

\end{v3topic}
\begin{v3topic}{Anonymisation vs.\ Pseudonymisation}
    \begin{tblr}{|X[2]|X[4]|X[4]|}
        \hline
        & \textbf{Anonymisation} & \textbf{Pseudonymisation} \\
        \hline
        Definition & Re-identification is \emph{irreversibly impossible} & Replaced with artificial identifier; re-identification possible with key \\
        \hline
        GDPR applicability & GDPR does \emph{not} apply & GDPR still applies; counts as a safeguard \\
        \hline
        ML use & Freely usable for training & Reduces risk but not a free pass \\
        \hline
        Challenge & True anonymisation is technically hard (re-identification attacks) & Key management is a critical security requirement \\
        \hline
    \end{tblr}
\captionof{table}{GDPR for AI/ML: Anonymisation vs.\ Pseudonymisation.}
\label{tab:ch17-it-law-inline-05}
    Re-identification of ``anonymised'' datasets via auxiliary information is a well-documented ML attack vector\textsuperscript{\cite[][pp.~1--15]{voightGDPRPracticalGuide2017}}.
\end{v3topic}

\subsection{International Data Transfers}

\begin{v3topic}{Transferring Data Outside the EEA}
    GDPR restricts transfers of personal data to countries outside the European Economic Area (EEA) unless adequate protection is ensured\textsuperscript{\cite[][pp.~200--230]{itgovernanceprivacyteamEUGDPR2020}}:
    \begin{tblr}{|X[2]|X[4]|}
        \hline
        \textbf{Mechanism} & \textbf{How it works} \\
        \hline
        Adequacy Decision & EU Commission declares destination country adequate (e.g.\ UK post-Brexit, Japan, Canada); no extra steps needed \\
        \hline
        Standard Contractual Clauses (SCCs) & EU-approved contract templates; most common mechanism for US/global transfers \\
        \hline
        Binding Corporate Rules (BCRs) & Internal codes for multinational groups; expensive to obtain \\
        \hline
        Derogations (Art.\ 49) & Explicit consent, vital interests, public interest; narrow scope \\
        \hline
    \end{tblr}
\captionof{table}{International Data Transfers: Transferring Data Outside the EEA.}
\label{tab:ch17-it-law-inline-06}
    \textbf{Schrems II (2020)}: The ECJ invalidated the EU-US Privacy Shield, requiring additional safeguards even when using SCCs.
    The EU-US Data Privacy Framework (2023) replaced it, but faces ongoing legal challenges.
    Cloud-based ML training that sends EU personal data to US servers must address this.
\end{v3topic}

% ---- 17.3 AI-Specific Regulation ---------------------------------------------------------------
\section{AI-Specific Regulation}

\subsection{EU AI Act}

\begin{v3topic}{EU AI Act: Risk-Based Classification}
    The \textbf{EU Artificial Intelligence Act} (Regulation 2024/1689) is a
    broad horizontal AI law with staged application dates.
    Prohibitions and AI-literacy duties began applying in February 2025;
    GPAI duties began in August 2025; enforcement and specified transparency
    rules began in August 2026. Under the 2026 timeline, Annex III high-risk
    rules apply from December 2027 and product-embedded high-risk rules from
    August 2028\textsuperscript{\cite{europeanCommissionAIAct2026}}.
    It classifies AI systems by risk level:
    \factvisual{
        \textbf{Classification is a legal analysis, not a model label.}
        The intended purpose, deployment context, actor role, affected people,
        and applicable definitions determine which duties attach.
        The same technical component can sit in different legal categories in
        different products or workflows.
    }{
        \begin{tikzpicture}[y=0.72cm]
            \path[fill=AlmanachRed!72,draw=white]
                (-1.3,3.3) -- (1.3,3.3) -- (0.8,4.1) -- (-0.8,4.1) -- cycle;
            \node[font=\scriptsize\sffamily,text=white] at (0,3.68)
                {unacceptable};
            \path[fill=AlmanachOchre!72,draw=white]
                (-2.0,2.2) -- (2.0,2.2) -- (1.3,3.2) -- (-1.3,3.2) -- cycle;
            \node[font=\scriptsize\sffamily,text=white] at (0,2.68)
                {high risk};
            \path[fill=AlmanachTeal!58,draw=white]
                (-2.8,1.1) -- (2.8,1.1) -- (2.0,2.1) -- (-2.0,2.1) -- cycle;
            \node[font=\scriptsize\sffamily,text=white] at (0,1.58)
                {transparency duties};
            \path[fill=AlmanachTeal!24,draw=white]
                (-3.6,0) -- (3.6,0) -- (2.8,1.0) -- (-2.8,1.0) -- cycle;
            \node[font=\scriptsize\sffamily] at (0,0.48)
                {minimal risk / voluntary practice};
            \draw[almanach arrow] (4.0,0.2) -- (4.0,3.8)
                node[midway,right,font=\tiny\sffamily,align=left]
                {obligations\\intensify};
        \end{tikzpicture}
        \visualcaption{A conceptual risk pyramid for the categories discussed
          here. It is an orientation aid, not a substitute for classification
          under the Act.}{fig:eu-ai-act-risk-pyramid}
    }
    \begin{tblr}{|X[2]|X[4]|X[4]|}
        \hline
        \textbf{Risk Level} & \textbf{Examples} & \textbf{Requirements} \\
        \hline
        \SetCell[r=1]{c} Unacceptable & Certain manipulative, exploitative,
        social-scoring, biometric, and predictive-policing practices &
        \textbf{Prohibited}, subject to the Act's definitions and exceptions \\
        \hline
        High Risk & HR/recruitment tools, credit scoring, medical devices, critical infrastructure, law enforcement, education assessment & Conformity assessment, human oversight, transparency, data governance, logging, incident reporting \\
        \hline
        Transparency Risk & Interactive AI and specified synthetic content &
        Disclosure or marking duties where Article 50 applies \\
        \hline
        Minimal Risk & Spam filters, AI-assisted games & No specific obligations; voluntary codes encouraged \\
        \hline
    \end{tblr}
\captionof{table}{EU AI Act: EU AI Act: Risk-Based Classification.}
\label{tab:ch17-it-law-inline-07}

\end{v3topic}
\begin{v3topic}{High-Risk AI: Key Obligations}
    Providers of high-risk AI systems must\textsuperscript{[EU AI Act, Arts.\ 9--17]}:
\visualseparator
    \begin{itemize}
        \item Establish a \textbf{risk management system} (ongoing)
        \item Implement \textbf{data governance}: training/validation data quality checks
        \item Produce \textbf{technical documentation} and keep logs
        \item Ensure \textbf{transparency} to users and regulators
        \item Enable \textbf{human oversight} mechanisms
        \item Achieve sufficient \textbf{accuracy, robustness, and cybersecurity}
        \item Register in the EU \textbf{AI database} before deployment
        \item Conduct \textbf{conformity assessment} (self-assessment for most; third-party for some)
    \end{itemize}

\end{v3topic}
\begin{v3topic}{General-Purpose AI (GPAI) Models}
    The AI Act also regulates \textbf{GPAI models} (e.g.\ large language models):
\visualseparator
    \begin{itemize}
        \item All GPAI providers: technical documentation, copyright compliance, training data summary
        \item GPAI with \textbf{systemic risk}: $10^{25}$ training FLOPs creates
              a statutory presumption, while other models can be designated;
              additional evaluation, incident, and cybersecurity duties apply
        \item Applies to models made available in the EU, regardless of where trained
        \item Open-source models have partial exemptions (transparency obligations remain)
    \end{itemize}
\end{v3topic}

\subsection{Other Jurisdictions}

\begin{v3topic}{Global AI Regulatory Landscape}
    \begin{tblr}{|X[2]|X[4]|}
        \hline
        \textbf{Jurisdiction} & \textbf{Approach} \\
        \hline
        \textbf{United States} & Federal sector and cross-cutting authorities
        coexist with state law. Executive Order 14179 revoked Executive Order
        14110 in January 2025 and directed a new federal AI policy review and
        action plan\textsuperscript{\cite{whiteHouseRemovingBarriersAI2025}}. \\
        \hline
        \textbf{China} & Multiple specific regulations: Deep Synthesis Provisions (deepfakes, 2022), Generative AI Regulation (2023), Algorithm Recommendation Regulation (2022). Focus on security, content control, data localisation. \\
        \hline
        \textbf{United Kingdom} & ``Pro-innovation'' approach: no single AI law; sector regulators (ICO, FCA, CMA) apply existing frameworks. AI Safety Institute established 2023. May diverge from EU. \\
        \hline
        \textbf{Brazil / Canada} & Privacy, consumer, sector, and evolving
        AI-specific proposals differ by jurisdiction; verify enacted text and
        effective dates before relying on a comparison. \\
        \hline
    \end{tblr}
\captionof{table}{Other Jurisdictions: Global AI Regulatory Landscape.}
\label{tab:ch17-it-law-inline-08}
\end{v3topic}

% ---- 17.4 Intellectual Property and AI ---------------------------------------------------------
\section{Intellectual Property and AI}

\subsection{Copyright and AI}

\begin{v3topic}{Training Data and Copyright}
    Scraping copyrighted text, images, or code to train AI models raises copyright infringement risks.
    Key questions\textsuperscript{\cite[][pp.~150--180]{lloydInformationTechnologyLaw2020}}:
\visualseparator
    \begin{itemize}
        \item \textbf{EU}: Text and Data Mining (TDM) exception (DSM Directive, Art.\ 4) allows mining of lawfully accessed content unless rights-holders opt out; Art.\ 3 allows TDM for research organisations
        \item \textbf{US}: Fair use doctrine — outcome uncertain; ongoing litigation (e.g.\ Getty Images v.\ Stability AI, NYT v.\ OpenAI)
        \item Best practice: use openly licensed datasets, document data provenance, respect opt-out mechanisms
    \end{itemize}

\end{v3topic}
\begin{v3topic}{AI-Generated Content and Authorship}
    Copyright requires a \textbf{human author} in most jurisdictions.
    Fully AI-generated output (no human creative input) is generally \emph{not} copyrightable in the EU or US\textsuperscript{\cite[][pp.~185--200]{lloydInformationTechnologyLaw2020}}.
\visualseparator
    \begin{itemize}
        \item EU: no authorship without a natural person; AI as tool can qualify if human made creative choices
        \item US Copyright Office: rejects purely AI-generated works; hybrid works assessed case-by-case
        \item UK: Computer-generated works have a narrow copyright (50 years, person who made arrangements)
        \item Implication: AI-generated code, images, text may be in the public domain
    \end{itemize}
\end{v3topic}

\subsection{Software and Open Source Licensing}

\begin{v3topic}{Open-Source Licences for AI Models}
    \begin{tblr}{|X[2]|X[2]|X[6]|}
        \hline
        \textbf{Licence} & \textbf{Type} & \textbf{Key Obligations / Restrictions} \\
        \hline
        MIT & Permissive & Attribution only; can use in proprietary products \\
        \hline
        Apache 2.0 & Permissive & Attribution + patent grant; popular for ML frameworks (TensorFlow) \\
        \hline
        GPL v3 & Copyleft & Derivatives must also be GPL; ``viral'' — problematic for proprietary AI products \\
        \hline
        LGPL & Weak copyleft & Library linking allowed without copyleft propagation \\
        \hline
        RAIL (Responsible AI Licence) & Use-restriction & Adds \emph{use restrictions} (e.g.\ no military, no surveillance) on top of permissive base; used by BLOOM, Stable Diffusion \\
        \hline
        Llama Community Licence & Proprietary & Usage restrictions above certain user thresholds; not OSI-approved \\
        \hline
    \end{tblr}
\captionof{table}{Software and Open Source Licensing: Open-Source Licences for AI Models.}
\label{tab:ch17-it-law-inline-09}
    \textsuperscript{\cite[][pp.~210--240]{lloydInformationTechnologyLaw2020}} Always check: (1) model licence, (2) training data licence, (3) code licence — they are independent.
\end{v3topic}

% ---- 17.5 Cybersecurity and Liability ----------------------------------------------------------
\section{Cybersecurity and Liability}

\subsection{Information Security Law}

\begin{v3topic}{NIS2 Directive (EU 2022/2555)}
    The \textbf{NIS2 Directive} (transposed by EU member states by Oct.\ 2024) significantly expands cybersecurity obligations for ``essential'' and ``important'' entities\textsuperscript{\cite[][pp.~280--310]{savinEUInternetLaw2017}}:
\visualseparator
    \begin{itemize}
        \item Covers: energy, transport, health, digital infrastructure, cloud providers, managed security services, and more
        \item Requires: risk management measures, supply-chain security, encryption, access control, incident response plans
        \item \textbf{Incident reporting}: significant incidents must be reported to national authority within 24 hours (early warning) and 72 hours (full notification)
        \item Senior management accountability: personal liability for non-compliance
        \item Fines: up to €10M or 2\% global turnover (essential entities)
    \end{itemize}

\end{v3topic}
\begin{v3topic}{eIDAS and Electronic Signatures}
    \textbf{eIDAS} (Regulation 910/2014, updated eIDAS2 under way) establishes a framework for electronic identification and trust services across the EU\textsuperscript{\cite[][pp.~320--340]{savinEUInternetLaw2017}}:
    \begin{tblr}{|X[2]|X[4]|}
        \hline
        \textbf{Type} & \textbf{Legal Status} \\
        \hline
        Simple electronic signature & Basic; low assurance \\
        \hline
        Advanced electronic signature (AdES) & Uniquely linked to signatory; tamper-evident \\
        \hline
        Qualified electronic signature (QES) & Equivalent to handwritten signature; highest assurance \\
        \hline
    \end{tblr}
\captionof{table}{Information Security Law: eIDAS and Electronic Signatures.}
\label{tab:ch17-it-law-inline-10}
    Relevant for AI systems that automate contract execution or identity verification.
\end{v3topic}

\subsection{AI Liability}

\begin{v3topic}{Liability Gaps in AI Systems}
    Traditional product liability (EU Product Liability Directive) requires proof of defect, damage, and causation.
    AI creates accountability gaps\textsuperscript{\cite[][pp.~350--370]{lloydInformationTechnologyLaw2020}}:
    \begin{tblr}{|X[2]|X[4]|}
        \hline
        \textbf{Challenge} & \textbf{Explanation} \\
        \hline
        Opacity & ML models may not provide traceable causal chains \\
        \hline
        Distributed responsibility & Data provider, model developer, deployer, and user all contribute \\
        \hline
        Autonomous behaviour & AI actions may not be foreseeable from design \\
        \hline
        Evolving systems & Models that update post-deployment may differ from what was tested \\
        \hline
    \end{tblr}
\captionof{table}{AI Liability: Liability Gaps in AI Systems.}
\label{tab:ch17-it-law-inline-11}
    \textbf{EU AI Liability Directive (proposed, 2022)}: Would introduce a rebuttable \emph{presumption of causality} — if a claimant can show an AI system violated a legal duty and caused harm, the burden shifts to the provider to disprove causation.
    Combined with the AI Act's obligations, non-compliant high-risk AI faces both regulatory fines \emph{and} civil liability exposure.
\end{v3topic}

% ---- 17.6 Further Reading -----------------------------------------------------------------------
\section{Literature and Further Reading}

\begin{v3topic}{Further Reading}
    \begin{itemize}
        \item \fullcite{lloydInformationTechnologyLaw2020} --- Comprehensive UK/EU IT law reference; covers data protection, IP, e-commerce, cybersecurity.
        \item \fullcite{itgovernanceprivacyteamEUGDPR2020} --- Practical GDPR implementation guide; articles by article walk-through.
        \item \fullcite{voightGDPRPracticalGuide2017} --- Legal practitioner's GDPR guide; good on data transfers and enforcement.
        \item \fullcite{savinEUInternetLaw2017} --- EU internet law: e-commerce, data protection, network regulation.
        \item \fullcite{lutziPrivateInternationalLaw2020} --- Jurisdiction and conflict-of-laws in online systems; relevant for cross-border AI services.
        \item \fullcite{siemsComparativeLaw2018} --- Comparative legal methods; useful for understanding divergence between EU/US/other regimes.
        \item \fullcite{thirlwaySourcesInternationalLaw2019} --- Foundations of international law; context for where EU law sits in the global system.
    \end{itemize}
\end{v3topic}

%
\section{Deep Dive: From Legal Classification to Engineering Evidence}
\begin{v3topic}{Begin with Role and Intended Purpose}
Map jurisdictions, actor roles, intended purpose, affected people, sector,
data flows, and lifecycle stage before selecting obligations.
Provider, deployer, importer, distributor, employer, controller, and processor
can carry different duties for the same technical component.

\end{v3topic}
\begin{v3topic}{Translate Duties into Controls}
Turn requirements into owners, controls, evidence, retention, review dates, and
incident paths.
Examples include provenance records, risk management, human oversight,
cybersecurity, performance monitoring, notices, rights handling, and supplier
contracts.

\end{v3topic}
\begin{v3topic}{Maintain a Legal Change Log}
Record the source, jurisdiction, version, effective date, interpretation owner,
affected system, and next review.
The legal chapter is orientation, not legal advice; current official text and
qualified counsel govern consequential decisions.
Reassess when purpose, data, model, geography, users, or autonomy changes.
\end{v3topic}

\chapterreview{Classify}{Identify}{Document}{Reassess}
